\documentclass[10pt,a4paper]{amsart}
\usepackage[margin=19mm]{geometry}
\usepackage{amsmath,amssymb,mathtools}
\usepackage[scr=rsfs,cal=euler]{mathalfa}
\usepackage{xfrac}
\usepackage[unicode,hidelinks,bookmarks=false]{hyperref}
\newcommand{\ie}{i.e.\ }
\newcommand{\cf}{cf.\ }
\newcommand{\resp}{respectively\ }
\newcommand{\Subsection}{\subsection}
\providecommand{\nobreakdash}{\nobreak\hbox{-}}
\setlength{\emergencystretch}{2em}
\multlinegap0pt
\usepackage{dsfont}
\usepackage{amssymb}
\usepackage{enumerate}
\usepackage{mathtools}
\usepackage[shortlabels]{enumitem}
\usepackage{xcolor}
\newtheorem{proposition}{Proposition}
\def\lm{{\lambda}}
\title{Linear stability of the elliptic relative equilibria for the restricted $N$-body problem: two special cases}
\author{Jiashengliang Xie, Bowen Liu, Qinglong Zhou}
\date{Source: arXiv:2310.00286v2 (CC BY 4.0)}
\begin{document}
\maketitle

\noindent\emph{Source and licence.} Linear stability of the elliptic relative equilibria for the restricted $N$-body problem: two special cases. Version arXiv:2310.00286v2 (CC BY 4.0), distributed by arXiv under the Creative Commons Attribution 4.0 International licence, http://creativecommons.org/licenses/by/4.0/, as recorded in the arXiv licence metadata for this exact version and shown on the abstract page https://arxiv.org/abs/2310.00286v2. Full licence notice: \texttt{licenses/arxiv-2310.00286-CC-BY-4.0.txt}.\par\smallskip

\noindent\textit{Editorial scope.} Counted unit: the article's proposition polygon.prop.lm3.lm4 (source0.tex lines 1255--1263). The article's complete original proof of it is delivered with it (source0.tex lines 1265--1336, one \texttt{proof} environment of 72 lines). No mathematics is written, repaired or re-derived: the statement and the proof are the article's own lines, in the article's own notation, and no retained line is substituted or re-flowed.  Retained uncounted as the source-local context the proof needs: lines 321--326.  Nothing was omitted from any retained span, so the package carries no omission record.  No retained line contains a citation, so the wrapper carries no bibliography and no bibliography entry.  No retained line contains a figure environment, an image-inclusion command or drawing code, so no image is delivered and the package declares no asset dependency.  Editorial formatting only, in addition: an AMS \texttt{amsart} preamble that restates the article's own declarations and macros used by the retained lines, namely the environments \texttt{proposition} on separate counters, together with the standard packages those lines need.  The retained environments print in this document as Proposition 1 in order of appearance, whereas the article prints the same items with the numbering of its own full document.  The complete native source of the article as distributed in the arXiv e-print for this exact version is bundled unchanged as \texttt{sources/147800/source0.tex}.
  \begin{align}
    \label{matrix.D}
  D=I_2-\frac{1}{\mu}\left(\sum_{i=1}^{N-1}
  \frac{m_i}{|a_i-a_{N}|^3}\right)I
  +\frac{3}{\mu}\sum_{i=1}^{N-1} m_i\frac{(a_i-a_{N})(a_i-a_{N})^T}{|a_i-a_{N}|^5}.
  \end{align}

 \begin{proposition}\label{polygon.prop.lm3.lm4}
	Suppose that $a_i$ for $1\leq i \leq N-1$ forms a $(1+n)$-gon configuration. We have
    $$\mu\alpha^3=\omega^2.$$
    Moreover, the eigenvalues $\lm_3$ and $\lm_4$ of $D$ defined by \eqref{matrix.D} satisfy
	\begin{align}
	\lm_3 = 1+{A\over\omega^2}+{B\over\omega^2},\quad
	\lm_4 = 1+{A\over\omega^2}-{B\over\omega^2}={l_3(x,\theta)\over\omega^2}.
	\end{align}
\end{proposition}

 \begin{proof}
First, we have
 
 \begin{displaymath}
 tr(D)=\sum_{i=1}^{N-1}\frac{m_i}{|a_i-a_N|^3}  
 =\frac{1}{\alpha^3}\sum_{i=1}^{N-1}\frac{m}{|\omega_i-\omega_N|^3}+\frac{1}{\alpha^3}\frac{m_0}{|\omega_N|^3}
 =\frac{1}{\alpha^3}\cdot 2A,
 \end{displaymath}
 and
 \begin{eqnarray*}
 	&&
 	\sum_{i=1}^{N-1}m_i\frac{(a_i-a_N)(a_i-a_N)^T}{|a_i-a_N|^5}
 	\\
 	&=&
 	\sum_{i=1}^{N-2}m\frac{\frac{1}{2}\alpha^2|\omega_i-\omega_N|^2I_2+\frac{1}{2}\alpha^2\psi((\omega_i-\omega_N)^2)}
 	{\alpha^5|\omega_i-\omega_ N|^5}+
 	m_0\frac{\frac{1}{2}\alpha^2|\omega_N|^2+\frac{1}{2}\alpha^2\psi(\omega_N^2)}{\alpha^5|\omega_N|^5}
 	\\
 	&=&
 	\sum_{i=1}^{N-2}m(\frac{1}{2\alpha^3|\omega_i-\omega_N|^3}I_2+\frac{\psi((\omega_i-\omega_N)^2)}{2\alpha^3|\omega_i-\omega_N|^5})+
 	m_0(\frac{1}{2\alpha^3|\omega_N|^3}I_2+\frac{\psi((\omega_N)^2)}{2\alpha^3|\omega_N|^5})\\
 	&=&
 	[\sum_{i=1}^{N-2}\frac{m}{2\alpha^3|\omega_i-\omega_j|^3}I_2+\frac{m_0}{2\alpha^3|\omega|^3}I_2]+
 	[\sum_{i=1}^{N-2}m\frac{\psi((\omega_i-\omega_N)^2)}{2\alpha^3|\omega_i-\omega_N|^5}+\frac{\psi(\omega_N^2)}{2\alpha^3|\omega_N|^5}]\\
 	&=&
 	\frac{1}{\alpha^3}AI_2+\frac{1}{3\alpha^3}\psi(B).
 \end{eqnarray*}
 Therefore, we have the following formula
 \begin{eqnarray*}
 	D&=&I_2-\frac{1}{\mu}(\sum_{i=1}^{N-1}\frac{m_i}{|a_i-a_N|^3})I_2+\frac{3}{\mu}(\sum_{i=1}^{N-1}m_i\frac{(a_i-a_N)(a_i-a_N)^T}{|a_i-a_N|^5})\\
 	&=&I_2-\frac{1}{\mu\alpha^3}\cdot 2AI_2+\frac{3}{\mu\alpha^3}(AI_2+\frac{1}{3}\psi(B))\\
 	&=&I_2+\frac{1}{\mu\alpha^3}(AI_2+\psi(B))\\
 	&=&I_2+\frac{A}{\mu\alpha^3}I_2+\frac{1}{\mu\alpha^3}\psi(B)\\
 	&=&I_2+\frac{1+\beta_{2.0}}{2}I_2+\psi(\beta_{22.0}).
 \end{eqnarray*}
 This implies:
 \begin{eqnarray*}
 	\frac{1+\beta_{2.0}}{2}=\frac{A}{\mu\alpha^3},\quad 
 	\beta_{22.0}=\frac{B}{\mu\alpha^3}.
 \end{eqnarray*}
 Let $u=\frac{1+\beta_{2.0}}{2}$ and $ v=\beta_{22.0}$, hence we obtain
 \begin{displaymath}
 uI_2+\psi(v)=\left(\begin{array}{cc}
 u+v_x&v_y\\
 v_y&u-v_x
 \end{array}
 \right),
 \end{displaymath}
 where $v_x=\text{Re}\,\beta_{22.0}$ and $ v_y=\text{Im}\,\beta_{22.0}$, indicating
 $\text{tr}(uI_2+\psi(v))=2u$ and $\text{det}(uI_2+\psi(v))=u^2-|v|^2$. 
 Therefore, the characteristic polynomial  of $uI_2+\psi(v)$ is $\lambda^2-2u\lambda+u^2-|v|^2$,
 and then the eigenvalues of the matrix $uI_2+\psi(v)$ are $u\pm|v|$. 
 Thus the eigenvalues of D are
 \begin{eqnarray*}
 	\lambda_{3,4}=1+\frac{1+\beta_{2.0}}{2}\pm|\beta_{22.0}|=1+\frac{A}{\mu\alpha^3}\pm\frac{|B|}{\mu\alpha^3}.
 \end{eqnarray*}
 The following formulas are the direct results from above arguments.
 \begin{eqnarray*}
 	\mu&=&\sum_{1\leq i<j\leq N-1}\frac{m_im_j}{|a_i-a_j|}\\
 	&=&\sum_{1\leq i<j\leq N-2}\frac{m^2}{\alpha|\omega_i-\omega_j|}+\sum_{i=1}^{N-2}\frac{mm_0}{\alpha\omega_i},
 \end{eqnarray*}
 then
 \begin{eqnarray*}
 	\mu\alpha^3&=&\alpha^2(\sum_{1\leq i<j\leq N-2}\frac{m^2}{\omega_i-\omega_j}+\sum_{i=1}^{N-2}\frac{mm_0}{|\omega_i|})
 	\\&=&\frac{m}{N-2}\sum_{1\leq i<j\leq N-2 }\frac{1}{|\omega_i-\omega_j|}+m_0
 	\\&=&\frac{m}{N-2}\sum_{i<j}\frac{1}{|\omega_i-\omega_j|}+m_0
 	\\&=&\frac{m}{2}\sum_{1<j\leq N-2}\frac{1}{|\omega_0-\omega_j|}+m_0
 	\\&=&\omega ^2.
 \end{eqnarray*}
 Consequently, $\lambda_4=\frac{l_3(x,\theta)}{\mu\alpha^3}
 =\frac{l_3(x,\theta)}{\omega^2}$, where $l_3(x,\theta)=A+\omega^2-|B|$.
\end{proof}

\end{document}
